\section{Introduction}
\label{sec:introduction}

Frequency estimation lies on the critical path of protection, disturbance monitoring, converter supervision, and fast frequency response in low-inertia grids. As inverter-based resources displace synchronous machines, the available response window after a disturbance shrinks and the measured waveform becomes less benign: the estimator may be forced to react to phase discontinuities, harmonics, out-of-band tones, oscillatory recovery, or compound inverter-driven events on the same time scale as the protection decision itself~\cite{Milano2018LowInertia,Pinheiro2025,Ponce2025,Bernal2025,Zoghby2024}. Under those conditions, an estimator that looks excellent on a nominal sinusoidal test can still accumulate large transient error, remain above relay thresholds for too long, or exceed an acceptable compute budget.

That mismatch between disturbance physics and simplified reporting is a benchmark-design problem as much as an estimator-design problem. Classical standards-style scenarios remain necessary because they define common test points, but they do not by themselves answer the engineering question that matters in deployment: which estimator should be trusted when the disturbance family, the protection objective, and the computational budget are all explicit. A flat leaderboard cannot answer that question because it collapses cases in which the dominant failure mode is latency, cases in which it is spectral pollution, and cases in which it is transient false-trip risk.

This paper addresses that gap through a disturbance-aware benchmark study. The benchmark evaluates \ArchivalScenarioCount{} scenarios and \ArchivalEstimatorCount{} estimators, producing \ArchivalPairCount{} tuned estimator-scenario configurations and \ArchivalMonteCarloRuns{} Monte Carlo realizations per configuration. The study is not proposing a new estimator. It is establishing a more defensible way to compare existing estimators under low-inertia IBR stress by keeping the signal family, metric contract, and deployment objective visible from per-scenario evaluation to final interpretation.

The paper makes four concrete contributions.
\begin{enumerate}[leftmargin=1.8em]
\item It defines a benchmark protocol with per-pair tuning, dual-rate waveform generation and acquisition, scenario-wise Monte Carlo aggregation, and family-wise reporting that prevents easy families from washing out harder ones.
\item It organizes the study through explicit estimator and scenario taxonomies, making clear which algorithmic assumptions are being tested against ramps, magnitude steps, phase jumps, spectral interference, noisy ringdown, and composite IBR stress.
\item It formalizes six headline metrics spanning accuracy, worst-case transient behavior, relay-oriented trip-risk, recovery, and computational cost, and it complements descriptive comparisons with family-level statistical evidence.
\item It converts the benchmark outcome into deployment guidance by separating protection-oriented, low-cost synchronization-oriented, and metering-oriented operating regions rather than claiming a universal best estimator.
\end{enumerate}

The central scientific claim is therefore sharper than in the shorter paper: disturbance family is part of the problem definition, not a post hoc qualifier. SOGI-PLL leads ramps, EKF leads magnitude steps and out-of-band interference, RA-EKF leads modulation, SOGI-FLL leads phase jumps, UKF leads ringdown, and PLL leads the composite IBR multi-event case on RMSE. Those winners do not coincide with the trip-risk leaders, especially under phase jumps, harmonics, and compound IBR stress. The benchmark therefore exposes a disturbance-dependent latency-robustness trade space rather than a universal champion.

The rest of the paper is organized as follows. Section~\ref{sec:related_work} positions the study against standards-oriented and estimator-oriented literature. Section~\ref{sec:benchmark_scope} defines the benchmark design and evaluation protocol. Section~\ref{sec:taxonomy} explains the estimator and scenario taxonomies. Section~\ref{sec:metrics_and_hypotheses} details the metric design and statistical analysis. Sections~\ref{sec:results} and~\ref{sec:discussion} present and interpret the main findings. Section~\ref{sec:limitations} states the scope limits of the analysis, and the appendices provide full estimator and scenario catalogs.
